\section{THPT Như Thanh, Thanh Hóa}

\debai{
Trong mặt phẳng với hệ tọa độ $Oxy$ cho đường tròn $(C): x^2+y^2-2x+4y-20=0$ và đường thẳng $d:3x+4y-20=0$. Chứng minh rằng $d$ tiếp xúc với $(C)$. Tam giác $ABC$ có đỉnh $A\in (C)$, hai đỉnh $B,C\in d$, trung điểm cạnh $AB\in (C)$. Tìm tọa độ các đỉnh tam giác $ABC$, biết trực tâm tam giác $ABC$ trùng với tâm đường tròn $(C)$ và $B$ có hoành độ dương.
}

\debai{
Giải phương trình: $4\sqrt{5x^3-6x^2+2}+4\sqrt{-10x^3+8x^2+7x-1}+x-13=0$.
}

\debai{
Cho các số $a,b,c\in R, a^2+b^2+c^2\neq 0$ và $2(4a^2+4b^2+c^2)=(2a+2b+c)^2$.\\
Tìm giá trị lớn nhất, nhỏ nhất của biểu thức $A=\dfrac{8a^3+8b^3+c^3}{(2a+2b+2c)(4ab+2bc+2ca)}$.
}

\section{THPT Chuyên Hạ Long}

\debai{
Trong mặt phẳng với hệ trục tọa độ $Oxy$, cho tam giác $ABC$ có $A(2;6),B(1;1),C(6;3)$.\\
1. Viết phương trình đường tròn ngoại tiếp tam giác $ABC$.\\
2. Tìm trên các cạnh $AB,BC,CA$ các điểm $K,H,I$ sao cho chu vi tam giác $KHI$ nhỏ nhất.
}

\debai{
Giải hệ phương trình \\ \cl{$\begin{cases} 3y\sqrt{2+x}+8\sqrt{2+x}=10y-3xy+12\\ 5y^3\sqrt{2-x}-8=6y^2+xy^3\sqrt{2-x} \end{cases}$.} 
}

\debai{
Chứng minh rằng: Với mọi $\Delta ABC$ ta đều có\\
\cl{$\left ( \sin \dfrac{A}{2}+\sin \dfrac{B}{2}+ \sin \dfrac{C}{2} \right )\left ( \cot \dfrac{A}{2}+\cot \dfrac{B}{2}+ \cot \dfrac{C}{2} \right )\ge \dfrac{9\sqrt{3}}{2}$.} 
\c
}

\section{THPT chuyên Vĩnh Phúc, Vĩnh Phúc - Khối AB}

\debai{
Trong mặt phẳng với hệ trục tọa độ $Oxy$, cho tam giác cân $ABC$ có đáy $BC$ nằm trên đường thẳng $d: 2x-5y+1=0$, cạnh $AB$ nằm trên đường thẳng $d':12x-y-23=0$. Viết phương trình đường thẳng $AC$ biết nó đi qua điểm $M(3;1)$.
}

\debai{
Giải hệ phương trình \\ \cl{$\begin{cases} \sqrt{5x^2+2xy+2y^2}+\sqrt{2x^2+2xy+5y^2}=3(x+y)\\ \sqrt{x+2y+1}+2\sqrt[3]{12x+7y+8}=2xy+x+5 \end{cases}$.}
}

\debai{
Cho ba số thực dương $a,b,c$ thỏa mãn $a^2+b^2+c^2=3$.\\
Tìm giá trị nhỏ nhất của biểu thức $S=8(a+b+c)+5\left( \dfrac{1}{a}+\dfrac{1}{b}+\dfrac{1}{c}\right)$.
}

\section{THPT chuyên Vĩnh Phúc, Vĩnh Phúc - Khối D}

\debai{
Trong mặt phẳng với hệ trục tọa độ $Oxy$, cho đường tròn $(T): x^2+y^2-x-9y+18=0$ và hai điểm $A(4;1),B(3;-1)$. Gọi $C,D$ là hai điểm thuộc $(T)$ sao cho $ABCD$ là một hình bình hành. Viết phương trình đường thẳng $CD$.
}

\debai{
Giải hệ phương trình \\ \cl{$\begin{cases} x^3-y^3+6x-3y=3x^2+4\\ x^2+6y+19=2\sqrt{3x+4}+3\sqrt{5y+14} \end{cases}$.} 
}

\debai{
Cho $a,b,c$ là các số thực dương thỏa mãn $a^2+b^2+c^2=1$. Chứng minh bất đẳng thức:\\ \centerline{$\left(\dfrac{1}{a}+\dfrac{1}{b}+\dfrac{1}{c}\right)-(a+b+c)\ge 2\sqrt{3}$.} 
}

\section{THPT Hồng Quang, Hải Dương}

\debai{
Trong mặt phẳng tọa độ $Oxy$, cho tam giác $ABC$ có đường cao và đường trung tuyến kẻ từ đỉnh $A$ lần lượt có phương trình là $x-3y=0$ và $x+5y=0$. Đỉnh $C$ nằm trên đường thawgnr $\Delta: x+y-2=0$ và có hoành độ dương. Tìm tọa độ các đỉnh của tam giác $ABC$, biết đường thẳng chứa trung tuyến kẻ từ $C$ đi qua điểm $E(-2,6)$.
}

\debai{Giải hệ phương trình \\ \cl{$\begin{cases} x-\dfrac{1}{(x+1)^2}=\dfrac{y}{x+1}-\dfrac{1+y}{y}\\ \sqrt{8y+9}=(x+1)\sqrt{y}+2 \end{cases}$.} 

}

\debai{
Cho các số dương $x,y,z$ thỏa mãn $x>y$ và $(x+z)(y+z)=1$.\\
Tìm giá trị nhỏ nhất của biểu thức
$P=\dfrac{1}{(x-y)^2}+\dfrac{4}{(x+z)^2}+\dfrac{4}{(y+z)^2}.$
}

\section{THPT Lương Thế Vinh, Hà Nội }

\debai{
Trong mặt phẳng tọa độ $Oxy$, cho hình chữ nhật $ABCD$ có diện tích bằng $15$. Đường thẳng $AB$ có phương trình $x-2y=0$. Trọng tâm của tam giác $BCD$ là điểm $G\left (\dfrac{16}{3};\dfrac{13}{3}\right)$. Tìm tọa độ bốn đỉnh của hình chữ nhật biết điểm $B$ có tung độ lớn hơn $3$.
}

\debai{
Giải hệ phương trình \\ \cl{$\begin{cases} 2x^3-3+2\sqrt{y^2+3y}=2x\sqrt{y}+y\\ x^2-\sqrt{y+3}+\sqrt{y}=0 \end{cases}$.} 
}

\debai{
Cho các số thực $a,b$ không âm và thỏa mãn: $3(a+b)+2(ab+1)\ge 5(a^2+b^2)$.\\
Tìm giá trị lớn nhất của biểu thức $T=3\sqrt{a+b}-3(a^2+b^2)+2(a+b)-ab$.}

\section{THPT Nông Cống, Hà Nội }

\debai{
Trong mặt phẳng với hệ tọa độ $Oxy$, cho đường tròn $(C)$ tâm $I$ $(x_I>0)$, $(C)$ đi qua điểm $A(-2;3)$ và tiếp xúc với đường thẳng $(d_1):x+y+4=0$ tại điểm $B$. $(C)$ cắt $(d_2):3x+4y-16=0$ tại $C$ và $D$ sao cho $ABCD$ là hình thang có hai đáy là $AD$ và $BC$, hai đường chéo $AC,BD$ vuông góc với nhau. Tìm tọa độ các điểm $B,C,D$.
}

\debai{
Giải hệ phương trình \\ \cl{$\begin{cases} \sqrt{x^2+xy+2y^2}+\sqrt{y^2+xy+2x^2}=2(x+y)\\ (8y-6)\sqrt{x-1}=\left(2+\sqrt{y-2}\right)\left(y+4\sqrt{x-2}+3\right)  \end{cases}$.} 
}

\debai{
Cho $x,y$ là các số thực không âm thỏa mãn:\\
\cl{$\sqrt{2x^2+3xy+4y^2}+\sqrt{2y^2+3xy+4x^2}-3(x+y)^2\le 0$.}
Tìm giá trị nhỏ nhất của:\\
\cl{$P=2(x^3+y^3)+2(x^2+y^2)-xy+\sqrt{x^2+1}+\sqrt{y^2+1}$.}
}

\section{THPT Thường Xuân 3, Thanh Hóa}
\debai{\\
1. Trong mặt phẳng với hệ tọa độ $Oxy$, cho tam giác $ABC$ có đỉnh $A(3;-4)$. Phương trình đường trung trực cạnh $BC$, đường trung tuyến xuất phát từ $C$ lần lượt là $x+y-1=0$ và $3x-y-9=0$. Tìm tọa độ các đỉnh $B,C$ của tam giác $ABC$.\\
2. Trong mặt phẳng với hệ tọa độ $Oxy$, cho đường tròn $(C)$ có phương trình $x^2+y^2+2x-4y-8=0$ và đường thẳng $(\Delta)$ có phương trình: $2x-3y-1=0$. Chứng minh rằng $(\Delta)$ luôn cắt $(C)$ tại hai điểm phân biệt $A,B$. Tìm tọa độ điểm $M$ trên đường tròn $(C)$ sao cho diện tích tam giác $ABM$ lớn nhất.
}
\debai{
Tìm tất cả các giá trị thực của tham số $m$ để hệ sau có nghiệm thực:\\ \cl{$\begin{cases} x^2+\dfrac{4x^2}{(x+2)^2}\ge 5\\ x^4+8x^2+16mx+16m^2+32m+16=0 \end{cases}$.} 
}
\debai{
Tìm giá trị lớn nhất, giá trị nhỏ nhất của biểu thức\\
\cl{$P=\dfrac{\sqrt{5-4x}-\sqrt{1+a}}{\sqrt{5-4a}+2\sqrt{1+a}+6}$}
trong đó $a$ là tham số thực và $-1\le a\le \dfrac{5}{4}.$
}
\section{THPT Tĩnh Gia II}
\debai{
Trong mặt phẳng với hệ tọa độ $Oxy$, cho đường tròn $(C): x^2+y^2=5$ tâm $O$, đường thẳng $(d):3x-y-2=0$. Tìm tọa độ các điểm $A,B$ trên $(d)$ sao cho $OA=\dfrac{\sqrt{10}}{5}$ và đoạn $OB$ cắt $(C)$ tại $K$ sao cho $KA=KB$.
}

\debai{
Giải hệ phương trình $\begin{cases} \sqrt{x^2+2x+5}-\sqrt{y^2-2y+5}=y-3x-3\\y^2-3y+3=x^2-x  \end{cases}$. 
}

\debai{
Cho các số thực dương $a,b,c$. Chứng minh rằng:\\
\cl{$\dfrac{\sqrt{a+b+c}+\sqrt{a}}{b+c}+\dfrac{\sqrt{a+b+c}+\sqrt{b}}{c+a}+\dfrac{\sqrt{a+b+c}+\sqrt{c}}{a+b}\ge \dfrac{9+3\sqrt{3}}{2\sqrt{a+b+c}}$.}
}

\section{THPT Triệu Sơn 3, Thanh Hóa}

\debai{Trong mặt phẳng với hệ tọa độ $Oxy$ cho hình chữ nhật $ABCD$ có điểm $H(1;2)$ là hình chiếu vuông góc của $A$ lên $BD$. Điểm $M\left(\dfrac{9}{2};3\right)$ là trung điểm của cạnh $BC$, phương trình đường trung tuyến kẻ từ $A$ của $\Delta ADH$ là $d:4x+y-4=0$. Viết phương trình cạnh $BC$.}

\debai{
Giải hệ phương trình $\begin{cases} x\sqrt{x^2+y}+y=\sqrt{x^4+x^3}+x\\ x+\sqrt{y}+\sqrt{x-1}+\sqrt{y(x-1)}=\dfrac{9}{2}  \end{cases}$. 
}

\debai{
Cho $a,b,c$ thuộc khoảng $(0;1)$ thỏa mãn $\left (\dfrac{1}{a}-1\right)\left (\dfrac{1}{b}-1\right)\left (\dfrac{1}{c}-1\right)=1$. Tìm GTNN của biểu thức $P=a^2+b^2+c^2$.
}

\section{Trung tâm dạy thêm văn hóa Lê Hồng Phong}

\debai{Trong mặt phẳng $Oxy$ cho hình thang $ABCD$ có đáy lớn $CD=3AB$, $C(-3;-3)$, trung điểm của $AD$ là $M(3;1)$. Tìm tọa độ đỉnh $B$ biết $S_{BCD}=18, AB=\sqrt{10}$ và đỉnh $D$ có hoành độ nguyên dương.}

\debai{
Giải hệ phương trình $\begin{cases} x-y\sqrt{2-x}+2y^2=2\\ 2\left(\sqrt{x+2}-4y\right)+8\sqrt{y}\sqrt{xy+2y}=34-15x \end{cases}$. 
}

\debai{
Cho $x,y$ là các số không âm thỏa $x^2+y^2=2$. Tìm giá trị lớn nhất và nhỏ nhất của:\\
\cl{$P=5(x^5+y^5)+x^2y^2\left(5\sqrt{2xy+2}-4xy+12\right)$.}
}

\section{THPT chuyên Vĩnh Phúc, Vĩnh Phúc lần 2}

\debai{
Trong mặt phẳng với hệ tọa độ $Oxy$, cho đường tròn $(C): x^2+y^2-2x-4y-4=0$ tâm $I$ và điểm $M(3;2)$. Viết phương trình đường thẳng $\Delta$ đi qua $M$, $\Delta$ cắt $(C)$ tại hai điểm phân biệt $A,B$ sao cho diện tích tam giác $IAB$ lớn nhất.
}

\debai{
Giải hệ phương trình $\begin{cases} x^4-2x=y^4-y\\ (x^2-y^2)^3=3   \end{cases}$. 
}

\debai{
Cho các số $a,b,c$ không âm sao cho tổng hai số bất kì đều dương. Chứng minh rằng:\\
\cl{$\sqrt{\dfrac{a}{b+c}}+\sqrt{\dfrac{b}{c+a}}+\sqrt{\dfrac{c}{a+b}}+\dfrac{9\sqrt{ab+bc+ca}}{a+b+c}\ge 6$.}
}

\section{THPT Đồng Lộc, Hà Tĩnh}

\debai{ Trong mặt phẳng với hệ tọa độ $Oxy$, cho hình thang cân $ABCD$ có diện tích bằng $\dfrac{45}{2}$, đáy lớn $CD$ nằm trên đường thẳng $x-3y-3=0$. Biết hai đường chéo $AC,BD$ vuông góc với nhau tại $I(2;3)$. Viết phương trình đường thẳng chứa cạnh $BC$, biết điểm $C$ có hoành độ dương.}

\debai{
Giải hệ phương trình $\begin{cases} y^3+6y^2+16y-3x+11=0\\ x^3+3x^2+x+3y+3=0   \end{cases}$. 
}

\debai{
Cho $0<a,b,c<\dfrac{1}{2}$ thỏa mãn $a+2b+3c=2$. Chứng minh rằng:\\
\cl{$\dfrac{1}{a(4b+6c-3)}+\dfrac{2}{b(3c+a-1)}+\dfrac{9}{c(2a+4b-1)}\ge 54$.}
}

\section{THPT Hậu Lộc 2, Thanh Hóa}

\debai{
Trong không gian tọa độ $Oxyz$ cho mặt phẳng $(P): x-y+z+2=0$ và điểm $A(1;-1;2)$. Tìm tọa độ điểm $A'$ đối xứng với điểm $A$ qua mặt phẳng $(P)$. Viết phương trình mặt cầu đường kính $AA'$.
}

\debai{
Giải hệ phương trình \\ \cl{$\begin{cases} (x+1)^2+y^2=2\left(1+\dfrac{1-y^2}{x}\right)\\ 4y^2=(y^2-x^3+3x-2)\left(\sqrt{2-x^2}+1\right)   \end{cases}$.}
}

\debai{
Cho các số thực dương $x,y,z$ thỏa mãn $xy\ge 1, z\ge 1$. Tìm giá trị nhỏ nhất của biểu thức:\\
\cl{$P=\dfrac{x}{y+1}+\dfrac{y}{x+1}+\dfrac{z^3+2}{3(xy+1)}$}
}

\section{THPT Nguyễn Trung Thiên}

\debai{
Trong mặt phẳng với hệ tọa độ $Oxy$, cho tam giác $ABC$ có trung điểm cạnh $BC$ là $M(3;-1)$. Điểm $E(-1;-3)$ nằm trên đường thẳng $\Delta$ chứa đường cao qua đỉnh $B$. Đường thẳng $AC$ qua $F(1;3)$. Tìm tọa độ các đỉnh của $\Delta ABC$ biết đường tròn ngoại tiếp $\Delta ABC$ có đường kính $AD$ với $D(4;-2)$.
}

\debai{
Giải phương trình: $(x+2)\left(\sqrt{x^2+4x+7}+1\right)+x\left(\sqrt{x^2+3}+1\right)=0$.
}

\debai{
Cho $x,y,z$ là ba số thực dương thỏa mãn $x^2+y^2+z^2=3.$ Tìm giá trị nhỏ nhất của biểu thức sau:\\
\cl{$P=\dfrac{x}{(y+z)^2}+\dfrac{y}{(x+z)^2}+\dfrac{z}{(x+y)^2}$.}
}
